% !PRINTED none
% Part-title (Abschnitt-opening) page: no running head / page number printed on this leaf.

\begin{center}
\gesperrt{Eleventh Chapter.}

\textbf{\gesperrt{The Points of Inflection of a Curve of the Third Order.}}

\noindent\rule{0.25\linewidth}{0.4pt}

\S.\,85.

\gesperrt{Ternary Forms and Algebraic Curves.}
\end{center}

In order to be able, in the further course of our exposition,
to present to the reader some of the manifold and interesting
applications of algebra to geometric problems, without referring
him to other means, the necessary geometric foundations for
these applications shall here be derived in brief. In doing so we
restrict ourselves to the geometry of the plane, and even there
we set aside everything that is not directly required for our
applications\textsuperscript{1)}.

We employ the homogeneous triangular coordinates, which are
defined by this, that the position of a point $x$ is determined by
its distances from the three sides of the coordinate triangle,
measured in three arbitrary units of measure (or multiplied by
three arbitrary constant factors). These three quantities are called
the \gesperrt{coordinates} of the point $x$, and are usually denoted
by $x_1, x_2, x_3$ or by $y_1, y_2, y_3$ etc.

Between these three coordinates of the point $x$ there subsists
a linear, non-homogeneous relation, which one obtains by setting
the area of the coordinate triangle $(1, 2, 3)$ equal to the
sum of the areas of the three triangles $(x, 2, 3)$, $(x, 3, 1)$, $(x, 1, 2)$.
A coordinate, and hence also the expression for the
area of the corresponding one of the three last-named tri-
%CONT>>

\noindent\rule{0.25\linewidth}{0.4pt}

\footnotesize
\noindent\textsuperscript{1)} Fuller details may be found in the textbooks of analytic
geometry, among which the work of \gesperrt{George Salmon}, ``Higher plane
curves'', German by \gesperrt{Fiedler}, deserves special mention.
\normalsize
